% GENERATED FILE. Edit canonical lessons, then run npm run generate:books.
% canonical-source-hash: fnv1a64:d61aafe0e296b4a8
% canonical-lessons: ML-C48-teacher, ML-C48-student, ML-C48-doctor, ML-C48-farmer, ML-C48-guest

\chapter{Who Someone Is}
\label{ch:ml-roles}
\hlchaptermodality{48}

\begin{chapteropening}
\textbf{By the end of this chapter:} \emph{I can say what someone's work is in Malayalam, and recognise the masculine ending that three of the five words share.}

Complete ML-C48-guest and say all five words in order.
\end{chapteropening}

\section[adhyāpakan]{\ml{അധ്യാപകൻ} (adhyāpakan) — a teacher}

\label{lesson:ML-C48-teacher}



\begin{quote}
Before the new run of words: what did \emph{viral} mean, and what did \emph{vayaṟŭ} mean?
\end{quote}

\subsection*{You'll want to know: \ml{അധ്യാപകൻ}}
\textbf{\ml{അധ്യാപകൻ}} (\emph{adhyāpakan}) — \textquotedblleft{}a teacher\textquotedblright{}.

Sanskrit \emph{adhyāpaka}, 'one who causes another to study', with a Malayalam ending fastened on: the chillu \ml{ൻ} that marks a masculine noun and that this book has already met at the end of \ml{മകൻ} and \ml{മുത്തച്ഛൻ}.

The Sanskrit word is itself built out of a going-toward: \emph{adhi} plus a root about approaching, so a teacher is the one who brings the student up to the text.

Kannada picked a different Sanskrit word for the same job, \ka{ಶಿಕ್ಷಕ} (\emph{śikṣaka}). Two sisters, two borrowings, one classroom — which is a fair picture of how these languages have handled Sanskrit generally: not one loan, but a choice among several.

\begin{grammarlens}[title={What you've built}]
The first of five words for who someone is.
\end{grammarlens}

\subsection*{Your turn}
Say these aloud:

\begin{itemize}
\raggedright
  \item \emph{adhyāpakan}
  \item it once more, slowly
  \item \emph{adhyāpakan}, then \emph{makan}, and hear the same chillu close both
\end{itemize}

\begin{tcolorbox}[breakable,colback=teal!4,colframe=teal!35!black,title={Before you move on}]
What does \ml{അധ്യാപകൻ} mean? (\textquotedblleft{}a teacher\textquotedblright{}.) Say it once more.
\end{tcolorbox}
\section[vidyārtthi]{\ml{വിദ്യാർത്ഥി} (vidyārtthi) — a student}

\label{lesson:ML-C48-student}



\begin{quote}
Before the new one: what did \emph{adhyāpakan} mean?
\end{quote}

\subsection*{You'll want to know: \ml{വിദ്യാർത്ഥി}}
\textbf{\ml{വിദ്യാർത്ഥി}} (\emph{vidyārtthi}) — \textquotedblleft{}a student\textquotedblright{}.

Two Sanskrit pieces, and both are legible once pointed at: \emph{vidyā}, 'knowledge', and \emph{arthin}, 'one who seeks'. A student is a knowledge-seeker, said outright.

Kannada has \ka{ವಿದ್ಯಾರ್ಥಿ} (\emph{vidyārthi}), the identical borrowing. Where the teacher-words parted company, the student-words did not.

There is a chillu here too, the \ml{ർ} in the middle, doing the same job as the \ml{ൽ} of the leg-word: a consonant with no vowel after it, given its own shape rather than a mark underneath.

\begin{grammarlens}[title={What you've built}]
Two: a teacher and a student.
\end{grammarlens}

\subsection*{Your turn}
Say these aloud:

\begin{itemize}
\raggedright
  \item \emph{vidyārtthi}
  \item it once more, slowly
  \item \emph{adhyāpakan}, then \emph{vidyārtthi}, the pair that need each other
\end{itemize}

\begin{tcolorbox}[breakable,colback=teal!4,colframe=teal!35!black,title={Before you move on}]
What does \ml{വിദ്യാർത്ഥി} mean? (\textquotedblleft{}a student\textquotedblright{}.) Say it once more.
\end{tcolorbox}
\section[vaidyan]{\ml{വൈദ്യൻ} (vaidyan) — a doctor}

\label{lesson:ML-C48-doctor}



\begin{quote}
Before the new one: what did \emph{vidyārtthi} mean?
\end{quote}

\subsection*{You'll want to know: \ml{വൈദ്യൻ}}
\textbf{\ml{വൈദ്യൻ}} (\emph{vaidyan}) — \textquotedblleft{}a doctor\textquotedblright{}.

The same \emph{vidyā} as the last word, stretched into an adjective: \emph{vaidya} is 'the one versed in the knowledge', and the knowledge meant is the medical one. Then the masculine chillu \ml{ൻ} again.

Kannada has \ka{ವೈದ್ಯ} (\emph{vaidya}) without that ending. Malayalam adds it; Kannada leaves the Sanskrit word standing as it came.

In a Kerala street today the word usually points at a practitioner of the traditional medicine, and the everyday word for a hospital doctor is the English one, worn down into Malayalam sounds. Both live side by side, which is worth saying rather than tidying away.

\begin{grammarlens}[title={What you've built}]
Three. A teacher, a student, a doctor.
\end{grammarlens}

\subsection*{Your turn}
Say these aloud:

\begin{itemize}
\raggedright
  \item \emph{vaidyan}
  \item it once more, slowly
  \item \emph{vidyārtthi}, then \emph{vaidyan}, and find the piece they share
\end{itemize}

\begin{tcolorbox}[breakable,colback=teal!4,colframe=teal!35!black,title={Before you move on}]
What does \ml{വൈദ്യൻ} mean? (\textquotedblleft{}a doctor\textquotedblright{}.) Say it once more.
\end{tcolorbox}
\section[karṣakan]{\ml{കർഷകൻ} (karṣakan) — a farmer}

\label{lesson:ML-C48-farmer}



\begin{quote}
Before the new one: what did \emph{vaidyan} mean?
\end{quote}

\subsection*{You'll want to know: \ml{കർഷകൻ}}
\textbf{\ml{കർഷകൻ}} (\emph{karṣakan}) — \textquotedblleft{}a farmer\textquotedblright{}.

Sanskrit \emph{karṣaka}, from a root about dragging or drawing — the plough pulled through the ground. The farmer is named for the furrow rather than for the crop.

The ending is the chillu \ml{ൻ} once more, the third of these role-words to take it, and the \ml{ർ} near the front is the same chillu as in the student-word. Two of them in one short word.

Kannada went elsewhere again and says \ka{ರೈತ} (\emph{raita}), a word that came into the language along a different road entirely.

\begin{grammarlens}[title={What you've built}]
Four. A teacher, a student, a doctor, a farmer.
\end{grammarlens}

\subsection*{Your turn}
Say these aloud:

\begin{itemize}
\raggedright
  \item \emph{karṣakan}
  \item it once more, slowly
  \item \emph{karṣakan}, then \emph{adhyāpakan}, and count the chillus in each
\end{itemize}

\begin{tcolorbox}[breakable,colback=teal!4,colframe=teal!35!black,title={Before you move on}]
What does \ml{കർഷകൻ} mean? (\textquotedblleft{}a farmer\textquotedblright{}.) Say it once more.
\end{tcolorbox}
\section[atithi]{\ml{അതിഥി} (atithi) — a guest}

\label{lesson:ML-C48-guest}



\begin{quote}
Before the new one: what did \emph{karṣakan} mean?
\end{quote}

\subsection*{You'll want to know: \ml{അതിഥി}}
\textbf{\ml{അതിഥി}} (\emph{atithi}) — \textquotedblleft{}a guest\textquotedblright{}.

Sanskrit \emph{atithi}, and Kannada \ka{ಅತಿಥಿ} is the identical word — the same borrowing arriving in both languages unchanged.

The traditional analysis inside Sanskrit takes it apart as \emph{a-} plus \emph{tithi}, a lunar day: the one who has no appointed day, who turns up unannounced. Whether the etymology is historically right is argued over; what matters is that Sanskrit itself understood the word that way, and built a whole ethic of hospitality on it.

This word takes no masculine ending. It sits on the other side of the chapter from the four before it.

\begin{grammarlens}[title={What you've built}]
Five, and the run is closed: a teacher, a student, a doctor, a farmer, a guest.
\end{grammarlens}

\subsection*{Your turn}
Say these aloud:

\begin{itemize}
\raggedright
  \item \emph{atithi}
  \item it once more, slowly
  \item all five in order — \emph{adhyāpakan}, \emph{vidyārtthi}, \emph{vaidyan}, \emph{karṣakan}, \emph{atithi}
\end{itemize}

\begin{tcolorbox}[breakable,colback=teal!4,colframe=teal!35!black,title={Before you move on}]
What does \ml{അതിഥി} mean? (\textquotedblleft{}a guest\textquotedblright{}.) Say it once more.
\end{tcolorbox}
